% MT059 / answer-zh — canonical editable content.

\begin{enumerate}[start=1]

\item 当Gaussian NB在类别之间共享方差时，两种模型都能在给定特征空间中产生线性决策边界。

\item 两者都可以用最大似然估计训练，但逻辑回归最大化标签的条件似然；NB则通过类别先验和类条件似然拟合联合模型。

\item NB是生成式模型；逻辑回归直接建模类别后验，是判别式模型。

\item 两者都能扩展到多分类，也都可以加入正则化或参数先验。

\item 对具有\(d\)个特征的二分类任务，逻辑回归有\(d+1\)个系数，其中包括截距。按原解采用的\textbf{单个共享标量方差}约定，NB有\(2d+1\)个类条件参数，即两个均值向量和一个方差。若还学习二分类的类别先验，就再增加一个自由参数，总数为\(2d+2\)。

\end{enumerate}


如果“共享方差”指的是\textbf{每个特征各有一个类间共享的方差}，那么NB连同可学习先验共有\(3d+1\)个自由参数。作答时应说明采用哪种模型和计数约定。
